% Part: first-order-logic
% Chapter: natural-deduction
% Section: proving-things-quant

\documentclass[../../../include/open-logic-section]{subfiles}

\begin{document}

\olfileid{fol}{ntd}{prq}

\olsection{\usetoken{P}{derivation} met kwantoren}

\begin{ex}
Bij het werken met kwantoren moeten we ervoor zorgen dat we de
eigenvariabelevoorwaarde niet schenden. Soms moeten we daarvoor de
volgorde aanpassen waarin we bepaalde inferenties uitvoeren. In het
algemeen helpt het om regels waarop de eigenvariabelevoorwaarde van
toepassing is als eerste te behandelen; in het voltooide bewijs staan
zij lager.

Laten we bekijken hoe we !!a{derivation} geven van de !!{formula}
$\lexists[x][\lnot !A(x)] \lif \lnot \lforall[x][!A(x)]$.
Zoals gebruikelijk beginnen we met:
\begin{prooftree}
\AxiomC{}
\UnaryInfC{$\lexists[x][\lnot !A(x)]\lif \lnot \lforall[x][!A(x)]$}
\end{prooftree}
Eerst schrijven we op wat er nodig is om die laatste stap met de regel
\Intro{\lif} te rechtvaardigen.
\begin{prooftree}
\AxiomC{$\Discharge{\lexists[x][\lnot !A(x)]}{1}$}
\DeduceC{$\lnot \lforall[x][!A(x)]$}
\DischargeRule{\Intro{\lif}}{1}
\UnaryInfC{$\lexists[x][\lnot !A(x)]\lif \lnot \lforall[x][!A(x)]$}
\end{prooftree}
Omdat er geen voor de hand liggende regel is die we op
$\lnot \lforall[x][!A(x)]$ kunnen toepassen, richten we de
!!{derivation} zo in dat we de regel \Elim{\lexists} kunnen gebruiken.
Daarbij moeten we op de eigenvariabelevoorwaarde letten en een
constante kiezen die niet voorkomt in $\lexists[x][!A(x)]$ of in een
aanname waarvan deze formule afhangt. (Er komen echter geen
!!{constant}s voor, zodat elke keuze geschikt is.)
\begin{prooftree}
\AxiomC{$\Discharge{\lexists[x][\lnot !A(x)]}{1}$}
\AxiomC{$\Discharge{\lnot !A(a)}{2}$}
\DeduceC{$\lnot \lforall[x][!A(x)]$}
\DischargeRule{\Elim{\lexists}}{2}
\BinaryInfC{$\lnot \lforall[x][!A(x)]$}
\DischargeRule{\Intro{\lif}}{1}
\UnaryInfC{$\lexists[x][\lnot !A(x)] \lif \lnot \lforall[x][!A(x)]$}
\end{prooftree}
Om $\lnot \lforall[x][!A(x)]$ af te leiden, proberen we de regel
\Intro{\lnot} te gebruiken. Daarvoor moeten we een tegenspraak
afleiden, eventueel met $\lforall[x][!A(x)]$ als extra aanname.
Uiteraard mag die tegenspraak mede berusten op de aanname
$\lnot !A(a)$, die door de \Elim{\lexists}-inferentie zal worden
opgeheven. We kunnen de afleiding als volgt opzetten:
\begin{prooftree}
\AxiomC{$\Discharge{\lexists[x][\lnot !A(x)]}{1}$}
\AxiomC{$\Discharge{\lnot !A(a)}{2}, \Discharge{\lforall[x][!A(x)]}{3}$}
\DeduceC{$\lfalse$}
\DischargeRule{\Intro{\lnot}}{3}
\UnaryInfC{$\lnot \lforall[x][!A(x)]$}
\DischargeRule{\Elim{\lexists}}{2}
\BinaryInfC{$\lnot \lforall[x][!A(x)]$}
\DischargeRule{\Intro{\lif}}{1}
\UnaryInfC{$\lexists[x][\lnot !A(x)]\lif \lnot \lforall[x][!A(x)]$}
\end{prooftree}
We lijken dicht bij een tegenspraak te zijn. De eenvoudigst toe te
passen regel is \Elim{\lforall}, waarvoor geen
eigenvariabelevoorwaarde geldt. Omdat we de universeel
gekwantificeerde~$x$ door elke gewenste term mogen vervangen, ligt het
voor de hand~$a$ te blijven gebruiken, zodat we een tegenspraak
kunnen bereiken.
\begin{prooftree}
\AxiomC{$\Discharge{\lexists[x][\lnot !A(x)]}{1}$}
\AxiomC{$\Discharge{\lnot !A(a)}{2}$}
\AxiomC{$\Discharge{\lforall[x][!A(x)]}{3}$}
\RightLabel{\Elim{\lforall}}
\UnaryInfC{$!A(a)$}
\RightLabel{\Elim{\lnot}}
\BinaryInfC{$\lfalse$}
\DischargeRule{\Intro{\lnot}}{3}
\UnaryInfC{$\lnot \lforall[x][!A(x)]$}
\DischargeRule{\Elim{\lexists}}{2}
\BinaryInfC{$\lnot \lforall[x][!A(x)]$}
\DischargeRule{\Intro{\lif}}{1}
\UnaryInfC{$\lexists[x][\lnot !A(x)]\lif \lnot \lforall[x][!A(x)]$}
\end{prooftree}

Vooral bij kwantoren is het van belang nu nogmaals te controleren dat
de eigenvariabelevoorwaarde niet is geschonden. De enige toegepaste
regel waarop die voorwaarde van toepassing is, was \Elim{\exists}.
Omdat de eigenvariabele~$a$ niet voorkomt in een aanname waarvan de
betreffende formule afhangt, is dit een correcte !!{derivation}.
\end{ex}

\begin{ex}
Soms leiden we !!a{formula} af uit andere !!{formula}s. In zulke
gevallen kunnen er niet-opgeheven aannamen zijn. Het is belangrijk
zowel onze aannamen als het uiteindelijke doel bij te houden.

Laten we bekijken hoe we !!a{derivation} geven van de !!{formula}
$\lexists[x][!C(x,b)]$ uit de aannamen $\lexists[x][(!A(x) 
\land !B(x))]$ en $\lforall[x][(!B(x) \lif !C(x,b))]$.
Zoals gebruikelijk schrijven we de conclusie onderaan.
\begin{prooftree}
\AxiomC{}
\UnaryInfC{$\lexists[x][!C(x,b)]$}
\end{prooftree}

We beschikken over twee premissen. Om de eerste te gebruiken, dat wil
zeggen om !!a{derivation} van $\lexists[x][!C(x, b)]$ uit
$\lexists[x][(!A(x) \land !B(x))]$ te zoeken, gebruiken we de regel
\Elim{\lexists}. Omdat daarop een eigenvariabelevoorwaarde van
toepassing is, passen we die regel als eerste toe. We verkrijgen:
\begin{prooftree}
\AxiomC{$\lexists[x][(!A(x) \land !B(x))]$}
\AxiomC{$\Discharge{!A(a) \land !B(a)}{1}$}
\DeduceC{$\lexists[x][!C(x,b)]$}
\DischargeRule{\Elim{\lexists}}{1}
\BinaryInfC{$\lexists[x][!C(x,b)]$}
\end{prooftree}
De twee aannamen waarmee we werken, hebben~$!B$ gemeen. Het kan nu
nuttig zijn \Elim{\land} toe te passen om~$!B(a)$ afzonderlijk te
verkrijgen.
\begin{prooftree}
\AxiomC{$\lexists[x][(!A(x) \land !B(x)])$}
\AxiomC{$\Discharge{!A(a) 
\land !B(a)}{1}$}
\RightLabel{\Elim{\land}}
\UnaryInfC{$!B(a)$}
\DeduceC{$\lexists[x][!C(x,b)]$}
\DischargeRule{\Elim{\lexists}}{1}
\BinaryInfC{$\lexists[x][!C(x,b)]$}
\end{prooftree}

De tweede aanname waarmee we werken is~$\lforall[x][(!B(x) \lif
  !C(x,b))]$. Omdat er geen eigenvariabelevoorwaarde geldt, kunnen we
met \Elim{\lforall} voor $x$ de !!{constant}~$a$ invullen en zo
$!B(a) \lif !C(a, b)$ verkrijgen. We hebben nu zowel
$!B(a) \lif !C(a,b)$ als $!B(a)$. De volgende stap is een
rechtstreekse toepassing van de regel \Elim{\lif}.
\begin{prooftree}
\AxiomC{$\lexists[x][(!A(x) \land !B(x))]$}
\AxiomC{$\lforall[x][(!B(x) \lif !C(x,b))]$}
\RightLabel{\Elim{\lforall}}
\UnaryInfC{$!B(a) \lif !C(a,b)$}
\AxiomC{$\Discharge{!A(a) 
\land !B(a)}{1}$}
\RightLabel{\Elim{\land}}
\UnaryInfC{$!B(a)$}
\RightLabel{\Elim{\lif}}
\BinaryInfC{$!C(a,b)$}
\DeduceC{$\lexists[x][!C(x,b)]$}
\DischargeRule{\Elim{\lexists}}{1}
\insertBetweenHyps{\hspace{-5em}}
\BinaryInfC{$\lexists[x][!C(x,b)]$}
\end{prooftree}
We zijn bijna klaar!{} Met één toepassing van \Intro{\lexists}
bereiken we ons doel.
\begin{prooftree}
\AxiomC{$\lexists[x][(!A(x) \land !B(x))]$}
\AxiomC{$\lforall[x][(!B(x) \lif !C(x,b))]$}
\RightLabel{\Elim{\lforall}}
\UnaryInfC{$!B(a) \lif !C(a,b)$}
\AxiomC{$\Discharge{!A(a) 
\land !B(a)}{1}$}
\RightLabel{\Elim{\land}}
\UnaryInfC{$!B(a)$}
\RightLabel{\Elim{\lif}}
\BinaryInfC{$!C(a,b)$}
\RightLabel{\Intro{\lexists}}
\UnaryInfC{$\lexists[x][!C(x,b)]$}
\DischargeRule{\Elim{\lexists}}{1}
\insertBetweenHyps{\hspace{-5em}}
\BinaryInfC{$\lexists[x][!C(x,b)]$}
\end{prooftree}
Omdat we er bij elke stap op hebben gelet dat de
eigenvariabelevoorwaarden niet werden geschonden, kunnen we erop
vertrouwen dat dit een correcte !!{derivation} is.
\end{ex}

\begin{ex}
Geef !!a{derivation} van de !!{formula}
$\lnot\lforall[x][!A(x)]$ uit de aannamen $\lforall[x][!A(x)] 
\lif \lexists[y][!B(y)]$ en $\lnot\lexists[y][!B(y)]$.
Zoals gebruikelijk schrijven we de beoogde !!{formula} onderaan.
\begin{prooftree}
\AxiomC{}
\UnaryInfC{$\lnot\lforall[x][!A(x)]$}
\end{prooftree}
De laatste regel van de !!{derivation} is een negatie. We proberen
daarom \Intro{\lnot} te gebruiken. Daarvoor moeten we bepalen hoe we
een tegenspraak kunnen !!{derive}.
\begin{prooftree}
\AxiomC{$\Discharge{\lforall[x][!A(x)]}{1}$}
\DeduceC{$\lfalse$}
\DischargeRule{\Intro{\lnot}}{1}
\UnaryInfC{$\lnot\lforall[x][!A(x)]$}
\end{prooftree}
Tot zover gaat het goed. We zouden \Elim{\lforall} kunnen gebruiken,
maar het is niet duidelijk of we daarmee ons doel bereiken. Laten we
in plaats daarvan een van onze aannamen gebruiken.
$\lforall[x][!A(x)] \lif \lexists[y][!B(y)]$ stelt ons samen met
$\lforall[x][!A(x)]$ in staat de regel \Elim{\lif} te gebruiken.
\begin{prooftree}
\AxiomC{$\lforall[x][!A(x)] \lif \lexists[y][!B(y)]$}
\AxiomC{$\Discharge{\lforall[x][!A(x)]}{1}$}
\RightLabel{\Elim{\lif}}
\BinaryInfC{$\lexists[y][!B(y)]$}
\DeduceC{$\lfalse$}
\DischargeRule{\Intro{\lnot}}{1}
\UnaryInfC{$\lnot\lforall[x][!A(x)]$}
\end{prooftree}
Er rest ons nog één aanname. Daarmee kunnen we kennelijk via
\Elim{\lnot} een tegenspraak bereiken.
\begin{prooftree}
\AxiomC{$\lnot\lexists[y][!B(y)]$}
\AxiomC{$\lforall[x][!A(x)] \lif \lexists[y][!B(y)]$}
\AxiomC{$\Discharge{\lforall[x][!A(x)]}{1}$}
\RightLabel{\Elim{\lif}}
\BinaryInfC{$\lexists[y][!B(y)]$}
\RightLabel{\Elim{\lnot}}
\BinaryInfC{$\lfalse$}
\DischargeRule{\Intro{\lnot}}{1}
\UnaryInfC{$\lnot\lforall[x][!A(x)]$}
\end{prooftree}
\end{ex}

\begin{prob}
Geef !!{derivation}s die het volgende aantonen:
\begin{enumerate}
\item $\Proves (\lforall[x][!A(x)] \land \lforall[y][!B(y)]) \lif
\lforall[z][(!A(z) \land !B(z))]$.
\item $\Proves (\lexists[x][!A(x)] \lor \lexists[y][!B(y)]) \lif
\lexists[z][(!A(z) \lor !B(z))]$.
\item $\lforall[x][(!A(x) \lif !B)] \Proves \lexists[y][!A(y)] \lif !B$.
\item $\lforall[x][\lnot !A(x)] \Proves \lnot\lexists[x][!A(x)]$.
\item $\Proves \lnot\lexists[x][!A(x)] \lif \lforall[x][\lnot !A(x)]$.
\item $\Proves \lnot\lexists[x][\lforall[y][((!A(x,y) \lif \lnot
!A(y,y)) \land (\lnot !A(y,y) \lif !A(x,y)))]]$.
\end{enumerate}
\end{prob}

\begin{prob}
Geef !!{derivation}s die het volgende aantonen:
\begin{enumerate}
\item $\Proves \lnot\lforall[x][!A(x)] \lif \lexists[x][\lnot!A(x)]$.
\item $(\lforall[x][!A(x)] \lif !B) \Proves \lexists[y][(!A(y) \lif !B)]$.
\item $\Proves \lexists[x][(!A(x) \lif \lforall[y][!A(y)])]$.
\end{enumerate}
(Hiervoor is steeds de regel $\FalseCl$ nodig.)
\end{prob}
\end{document}
